%%%%%%%%%%%%%%%%%%%%%%%%%%%%%%%%%%%%%%%%%%%%%%%%%%%%%%%%%%%%%%%%%%%%%%
% Lớp 10 | Đại số tổ hợp | Quy tắc cộng-quy tắc nhân
% Dạng 0D8?1-1: Bài toán chỉ sử dụng quy tắc cộng
% Số câu: 20  (sinh tự động bởi tach_cau.py; từ nay thầy sửa trực tiếp file này)
%%%%%%%%%%%%%%%%%%%%%%%%%%%%%%%%%%%%%%%%%%%%%%%%%%%%%%%%%%%%%%%%%%%%%%

%%%%%%%%%%%%%%%% TRẮC NGHIỆM 4 PHƯƠNG ÁN %%%%%%%%%%%%%%%%

% nguồn: MaHoaNH_TracNghiem_Lop10
\begin{ex}%[0D8N1-1]%[ma:6d0a6672]
	Có $3$ cây bút đỏ, $4$ cây bút xanh trong một hộp bút. Hỏi có bao  nhiêu cách lấy ra $1$ cây bút từ hộp bút?
	\choice 
		{\True $7$} 
		 {$12$}
		 {$4$}
		 {$3$}
	 \loigiai{Hộp bút có $7$ cây bút nên có $7$ cách lấy $1$ cây bút từ hộp bút.}
\end{ex}

% nguồn: MaHoaNH_TracNghiem_Lop10
\begin{ex}%[0D8N1-1]%[ma:bafc22a8]
	Một lớp học có $15$ học sinh nam và $17$ học sinh nữ. Có bao nhiêu cách chọn ra một học sinh từ lớp học này để tham dự Đại hội Đoàn trường? 
	\choice 
		{$17$} 
		{$255$} 
		 {$15$}
		{\True $32$} 
	\loigiai{ 
		Số học sinh của lớp bằng $15+17=32$, chọn $1$ học sinh tham dự đại  hội đoàn trường có $32$ cách.
	 }
\end{ex}

% nguồn: MaHoaNH_TracNghiem_Lop10
\begin{ex}%[0D8N1-1]%[ma:c8f14a32]
	Lớp  10A có $25$ học sinh nam và $15$ học sinh nữ. Biết rằng tất cả các bạn trong lớp đều có khả năng tham gia vào đội thanh niên tình nguyện của trường. Để chọn ra một học sinh tham gia đội tình nguyện đó, ta tính số cách chọn bằng
	\choice 
		{\True $25+15$} 
		 {$25\cdot 15$}
		{$25-15$} 
		 {$25:15$}
	\loigiai{ 
		Để chọn một học sinh tham gia đội tình nguyện, ta  có thể chọn học sinh nam hoặc học sinh nữ, suy ra số cách chọn là $25+15=40$ cách.
	 }
\end{ex}

% nguồn: MaHoaNH_TracNghiem_Lop10
\begin{ex}%[0D8N1-1]%[ma:45be172a]
	Bạn An có $4$ quyển sách Toán học, $6$ quyển sách Văn  học và $5$ quyển sách Tiếng Anh, các quyển sách là khác nhau. Hỏi bạn An có bao nhiêu cách chọn một quyển sách để đọc?
	 \choice
		{$120$} 
		{$11$} 
		 {$10$}
		{\True $15$} 
	\loigiai{ 
	Bạn An có  số cách chọn một quyển sách để đọc là $4+6+5=15$.
	 }
\end{ex}

% nguồn: MaHoaNH_TracNghiem_Lop10
\begin{ex}%[0D8N1-1]%[ma:c266c9ab]
	Lớp 12A có $43$ học sinh, lớp 12B có $30$ học sinh. Hỏi có  bao nhiêu cách chọn $1$ học sinh từ lớp $12$A và $12$B?
	\choice 
		{$43$} 
		{$30$} 
		 {\True $73$}
		{$1290$} 
	 \loigiai{
		Mỗi cách chọn $1$ học sinh từ lớp $12$A và $12$B là một tổ hợp chập $1$ trong $43 + 30 = 73$ phần tử, do đó số cách chọn  là $\mathrm{C}_{73}^{1} = 73$.
	 }
\end{ex}

% nguồn: MaHoaNH_TracNghiem_Lop10
\begin{ex}%[0D8N1-1]%[ma:8f27040e]
	Một người vào cửa hàng ăn uống, người đó chọn một thực đơn gồm liên tiếp: $1$ món ăn trong $5$ món  khác nhau và $1$ nước uống trong $3$ loại nước khác nhau. Số cách lựa chọn thực đơn của người đó là
	\choice 
		{$8$} 
		 {$2$}
		{\True  $15$}
		 {$1$}
	 \loigiai{
		Chọn $1$ món ăn trong $5$ món khác  nhau có $5$ cách.\\
		Chọn $1$ nước uống trong $3$ loại nước khác  nhau có $3$ cách.\\
		Theo quy tắc nhân có tất cả $5\cdot  3=15$ cách chọn một thực đơn.
	} 
\end{ex}

% nguồn: MaHoaNH_TracNghiem_Lop10
\begin{ex}%[0D8N1-1]%[ma:eb446368]
	Bạn A  đi mua một cái áo thun ở một cửa hàng. Giả sử cửa hàng có $10$ cái áo thun size S, $5$ cái áo size M, $18$ áo size L. Giả sử áo size nào bạn A cũng mang được. Hỏi bạn A có bao nhiêu cách lựa chọn?
	\choice 
		{\True $33$  áo}
		{$900$ áo} 
		 {$15$ áo}
		{$23$ áo} 
	\loigiai{ 
		Vì bạn A có thể mang được  cả $3$ size nên bạn A chọn
		\begin{itemize}
			\item size S  có $10$ cách chọn,
			\item size M  có $5$ cách chọn,
			\item size L  có $18$ cách chọn.
		\end{itemize}
		Vậy theo quy tắc cộng, bạn A có $10 + 5 + 18 = 33$ cách lựa chọn. 
	 }
\end{ex}

% nguồn: MaHoaNH_TracNghiem_Lop10
\begin{ex}%[0D8H1-1]%[ma:9c7ac924]
	Trong một trường THPT, khối $11$ có $280$ học sinh nam và $325$ học sinh nữ. Nhà trường cần chọn một học sinh ở khối $11$  đi dự dạ hội của học sinh thành phố. Hỏi nhà trường có bao nhiêu cách chọn?
	\choice 
		{\True $605$} 
		 {$325$}
		{$280$} 
		 {$45$}
	\loigiai{ 
		Số học sinh khối $11$ của  trường là $280+325=605$.\\
		Vậy nhà trường có $605$ cách chọn một học sinh ở khối $11$ đi dự dạ hội của  học sinh thành phố.
	 }
\end{ex}

% nguồn: MaHoaNH_TracNghiem_Lop10
\begin{ex}%[0D8V1-1]%[ma:2caac5aa]
	Gọi $S$ là tập hợp tất cả các số tự nhiên gồm 5 chữ số đôi một khác nhau được lập từ các chữ số $5$, $6$, $7$, $8$, $9$. Tính tổng tất cả các  số thuộc tâp $S$.
	\choice 
		{$9\,333\,420$} 
		{$46\,666\,200$} 
		 {\True $9\,333\,240$}
		{$46\,666\,240$} 
	\loigiai{ 
		Số các số tự nhiên  gồm $5$ chữ số đôi một khác nhau được lập từ $5$, $6$, $7$, $8$, $9$ là $5!=120$ số.\\
		Vì vai trò các chữ số  như nhau nên mỗi chữ số $5$, $6$, $7$, $8$, $9$ xuất hiện ở hàng đơn vị là $4!=24$ lần.\\
		Tổng các chữ số ở hàng đơn vị  là $24\cdot(5+6+7+8+9)=840$.\\
		Tương tự thì mỗi  lần xuất hiện ở các hàng chục, trăm, nghìn, chục nghìn của mỗi chữ số là $24$ lần.\\
		Vậy tổng các  số thuộc tập $S$ là $840\cdot\left(1+10+10^2+10^3+10^4\right)=9\,333\,240$.
	} 
\end{ex}


%%%%%%%%%%%%%%%% ĐÚNG/SAI %%%%%%%%%%%%%%%%

% nguồn: MaHoaNH_DungSai_Lop10
\begin{ex}%[0D8N1-1]%[ma:c2ff79ed]
	Một thùng chứa $10$ quả cam, một thùng khác  chứa $7$ quả bưởi.
	 \choiceTF
		{\True Có  $17$ cách chọn một quả cam hoặc một quả bưởi}
		{\True Hai thùng  có tổng cộng $17$ quả}
		 {Có $70$ cách chọn một quả cam hoặc một quả bưởi}
		{Có $3$ cách chọn một quả cam hoặc một quả  bưởi}
	 \loigiai{
		Một thùng chứa $10$ quả cam,  một thùng khác chứa $7$ quả bưởi.\\
		Do đó hai thùng có tổng  cộng $17$ quả và có $17$ cách chọn một quả cam hoặc một quả bưởi.
	} 
\end{ex}

% nguồn: MaHoaNH_DungSai_Lop10
\begin{ex}%[0D8N1-1]%[ma:17de0f40]
	Trên bàn có $8$ chiếc bút chì khác nhau, $6$ chiếc bút bi khác nhau và $10$ cuốn tập khác  nhau. Khi đó, \textbf{đúng}, \textbf{sai} của các mệnh đề sau
	\choiceTF 
		{Có $480$ cách chọn một đồ vật duy nhất hoặc một cây bút  chì hoặc một cây bút bi hoặc một cuốn tập}
		{\True Nếu  chọn một cây bút chì thì có $8$ cách}
		{\True Nếu  chọn một cuốn tập thì có $10$ cách}
		{\True Nếu  chọn một cây bút bi thì có $6$ cách}
	 \loigiai{
		\begin{itemchoice}
			\itemch \textbf{Sai}. Theo quy tắc cộng, ta có  $8+6+10=24$ cách chọn.
			\itemch \textbf{Đúng}. Nếu chọn một  cây bút chì có $8$ cách.
			\itemch \textbf{Đúng}. Nếu chọn  một cuốn tập có $10$ cách.
			\itemch \textbf{Đúng}. Nếu chọn một cây  bút bi có $6$ cách.
		\end{itemchoice}
	 }
\end{ex}

% nguồn: MaHoaNH_DungSai_Lop10
\begin{ex}%[0D8N1-1]%[ma:bb78eb8d]
	Trong một cuộc thi tìm hiểu về đất nước Việt Nam, ban tổ chức công bố danh sách các đề tài bao gồm:  $6$ đề tài về lịch sử, $5$ đề tài về thiên nhiên, $4$ đề tài về con người và $3$ đề tài về văn hóa. Mỗi thí sinh được quyền chọn một đề tài. Khi đó, xét tính đúng sai của các mệnh đề sau
	\choiceTF 
		{\True Chọn  một đề tài thiên nhiên có $5$ cách}
		{Mỗi thí sinh có $360$ khả năng lựa chọn đề  tài}
		{\True Chọn một  đề tài văn hóa có $3$ cách}
		{Chọn  một đề tài lịch sử có $720$ cách}
	 \loigiai{
		\begin{itemchoice}
			\itemch \textbf{Đúng}. Chọn một đề tài thiên nhiên có $5$ cách. 
			\itemch \textbf{Sai}. Theo quy tắc cộng, số cách chọn một đề  tài là $6+5+4+3=18$ cách.
			\itemch \textbf{Đúng}. Chọn một  đề tài văn hóa có $3$ cách.
			\itemch \textbf{Sai}. Chọn một đề tài lịch sử  có $6$ cách.
		\end{itemchoice}
	} 
\end{ex}

% nguồn: MaHoaNH_DungSai_Lop10
\begin{ex}%[0D8H1-1]%[ma:4cc833f1]
	Hồng có $6$ quyển sách tham khảo môn Toán, $5$ quyển sách tham khảo môn  Văn, $3$ quyển sách tham khảo môn Tiếng Anh.
	\choiceTF 
		{\True Có $14$ cách chọn một quyển sách để quyên góp từ  thiện}
		{\True Có $11$ cách chọn một quyển sách tham khảo thuộc môn  Toán hoặc môn Văn}
		{\True Có $9$ cách chọn một quyển sách  tham khảo thuộc môn Toán hoặc môn Anh}
		{Có $90$ cách chọn một quyển sách tham khảo thuộc các môn  Toán, Văn, Anh}
	 \loigiai{
		Hồng có  $6$ quyển sách tham khảo môn Toán, $5$ quyển sách tham khảo môn Văn, $3$ quyển sách tham khảo môn Tiếng Anh nên
		\begin{itemize}
			\item Có $14$ cách chọn một quyển sách để quyên góp  từ thiện.
			\item  Có $11$ cách chọn một quyển sách tham khảo thuộc môn Toán hoặc môn Văn.
			\item Có $9$ cách chọn một quyển sách tham khảo thuộc môn  Toán hoặc môn Anh.
		\end{itemize}
	} 
\end{ex}


%%%%%%%%%%%%%%%% TRẢ LỜI NGẮN %%%%%%%%%%%%%%%%

% nguồn: MaHoaNH_TraLoiNgan_Lop10
\begin{bt}%[0D8N1-1]%[ma:5b031117]
	Trong một trường THPT, khối $11$ có $280$ học sinh nam và $325$ học sinh nữ. Nhà trường cần chọn một học sinh ở khối $11$ đi dự dạ hội của học sinh thành phố. Hỏi nhà trường có bao  nhiêu cách chọn?
	\shortans{605} 
	 \loigiai{
		Phương án  $1$: Nếu chọn một học sinh nam có $280$ cách.\\
		Phương án  $2$: Nếu chọn một học sinh nữ có $325$ cách.\\
		Theo quy tắc cộng, ta có  $280+325=605$ cách chọn.
	 }
\end{bt}

% nguồn: MaHoaNH_TraLoiNgan_Lop10
\begin{bt}%[0D8N1-1]%[ma:8b8cb7ac]
	Trong một cuộc thi tìm hiểu về đất nước  Việt Nam, ban tổ chức công bố danh sách các đề tài bao gồm: $8$ đề tài về lịch sử, $7$ đề tài về thiên nhiên, $10$ đề tài về con người và $6$ đề tài về văn hóa. Mỗi thí sinh được quyền chọn một đề tài. Hỏi mỗi thí sinh có bao nhiêu khả năng lựa chọn đề tài?
	 \shortans{31}
	 \loigiai{
		\begin{enumerate}[TH1:]
			\item Nếu chọn đề tài về lịch  sử có $8$ cách.
			\item Nếu  chọn đề tài về thiên nhiên có $7$ cách.
			\item Nếu chọn đề tài  về con người có $10$ cách.
			\item Nếu chọn đề tài  về văn hóa có $6$ cách.
		\end{enumerate}
		Theo quy tắc cộng,  ta có $8+7+10+6=31$ cách chọn.
	 }
\end{bt}


%%%%%%%%%%%%%%%% TỰ LUẬN %%%%%%%%%%%%%%%%

% nguồn: MaHoaNH_TuLuan_Lop10
\begin{bt}%[0D8H1-1]%[ma:b881360d]
	Tính số giao điểm tối đa của $10$ đường thẳng  phân biệt khi không có ba đường nào đồng quy và hai đường nào song song?
	\loigiai{ 
		\begin{itemize}
			\item Đường thẳng thứ nhất giao với  $9$ đường còn lại nên có $9$ giao điểm.
			\item Đường thẳng thứ hai  giao với $8$ đường còn lại nên có thêm $8$ giao điểm (đã tính giao điểm với đường thẳng thứ nhất ở trên).
			\item Đường thẳng thứ $9$ giao với $1$ đường còn lại nên có thêm  $1$ giao điểm.
		\end{itemize}
		 Vậy có $9+8+7+6+5+4+3+2+1=45$ giao điểm.
	} 
\end{bt}

% nguồn: MaHoaNH_TuLuan_Lop10
\begin{bt}%[0D8B1-1]%[ma:8edd42f7]
	Tung đồng thời hai con xúc xắc khác nhau và ghi lại số chấm xuất hiện trên mỗi con xúc xắc. Có bao nhiêu kết quả có thể xảy ra mà tổng số chấm xuất hiện trên hai mặt là bội  của 5?
	 \loigiai{
		Ta viết $(a,b)$ để kí hiệu kết quả số chấm xuất hiện trên hai con xúc xắc lần lượt là $a$ và $b$. Ta có $2 \leq a+b \leq 12$ nên $a+b$ là bội của 5 khi $a+b=5$ hoặc $a+b=10$.  Trường hợp $a+b=5$ gồm các kết quả: $(1,4),(2,3),(3,2),(4,1)$. Trường hợp này có 4 kết quả.\\
		Trường hợp $a+b=10$ bao gồm các kết quả: $(4,6),(5,5),(6,4)$. Trường hợp này có 3 kết  quả.\\
	Vậy có $4+3=7$ kết quả có thể xảy ra mà  tổng số chấm xuất hiện trên hai con xúc xắc là bội của 5.}
\end{bt}

% nguồn: MaHoaNH_TuLuan_Lop10
\begin{bt}%[0D8Y1-1]%[ma:0301baf1]
	Một nhóm tình nguyện viên gồm $4$ học sinh lớp $10$A, $5$ học sinh lớp  $10$B và $6$ học sinh lớp $10$C. Để tham gia một công việc tình nguyện, nhóm có bao nhiêu cách cử ra $1$ thành viên của nhóm?
	\loigiai{ 
		\begin{itemize}
			\item Số  cách cử $1$ thành viên trong $4$ học sinh lớp $10$A là $4$ cách.
			\item Số cách cử $1$ thành viên trong  $5$ học sinh lớp $10$B là $5$ cách.
			\item Số cách cử $1$ thành viên trong $6$ học  sinh lớp $10$C là $6$ cách.
		\end{itemize}
		Theo quy tắc cộng, ta có  $4+5+6=15$ cách.
	} 
\end{bt}

% nguồn: MaHoaNH_TuLuan_Lop10
\begin{bt}%[0D8Y1-1]%[ma:616dd6bd]
	Trong một cái hộp có chứa 8 quả bóng màu trắng đánh số từ 1 đến 8; 10 quả bóng màu xanh đánh số từ 1 đến 10; 12 quả bóng màu cam đánh số từ 1 đến 12. Từ hộp này,  có bao nhiêu cách
	\begin{enumerate}
		\item chọn ra một quả  bóng?
		\item  chọn ra ba quả bóng có màu khác nhau đôi một?
		\item chọn ra hai quả bóng có  màu khác nhau?
	\end{enumerate}
	\loigiai{ 
		\begin{enumerate}
			\item $8+10+12=30$. 
			 \item $8 \cdot 10 \cdot 12=960$.
			\item $8 \cdot 10+8 \cdot  12+10 \cdot 12=296$.
		\end{enumerate}
	 }
\end{bt}

% nguồn: MaHoaNH_TuLuan_Lop10
\begin{bt}%[0D8Y1-1]%[ma:1201f215]
	Cho tập $A=\{1;2;3;4;5;6;7\}$. Hỏi có bao nhiêu số tự nhiên có $4$ chữ số khác nhau  đôi một được lập từ $A$?
	\loigiai{ 
		\begin{itemize}
			\item Cách $1$: Dùng phép  đếm.\\
			Gọi $\overline{abcd}$  là số tự nhiên có $4$ chữ số khác nhau đôi một được lập từ $A$.
			\begin{itemize}
				\item $a\in A$ suy ra $a$ có $7$ cách  chọn.
				\item $b\in A\setminus \{a\}$ suy ra  $b$ có $6$ cách chọn.
				\item $c\in  A\setminus \{a,b\}$ suy ra $c$ có $5$ cách chọn.
				\item $d\in A\setminus \{a,b,c\}$ suy ra $d$  có $4$ cách chọn.
			\end{itemize}
			Theo quy tắc nhân, ta được $7\cdot6\cdot5\cdot4=840$ số. 
			\item Cách $2$: Dùng  chỉnh hợp.\\
			Số tự nhiên có $4$ chữ số khác nhau đôi  một được lập từ $A$ là một chỉnh hợp chập $4$ của $7$ phần tử.\\
			Vậy có $\mathrm{A}_7^4=7\cdot6\cdot5\cdot4=840$  số.
		\end{itemize}
	 }
\end{bt}
